% Final theory baseline: Sections II--IV.
% Compile with pdfLaTeX (twice for cross-references); IEEEtran supplies the text fonts.
% Compatibility layer for the actual Tectonic/XeTeX build only.
% pdfLaTeX ignores every conditional definition below.
% Define the conditional outside the XeTeX-only branch so TeX can skip it safely.
\newif\ifTKTermes
\ifdefined\XeTeXversion
  \TKTermestrue
  \IfFileExists{texgyretermes-regular.otf}{}{\TKTermesfalse}
  \IfFileExists{texgyretermes-italic.otf}{}{\TKTermesfalse}
  \IfFileExists{texgyretermes-bold.otf}{}{\TKTermesfalse}
  \IfFileExists{texgyretermes-bolditalic.otf}{}{\TKTermesfalse}
  \ifTKTermes
    \def\TKRomanRegular{[texgyretermes-regular.otf]}
    \def\TKRomanItalic{[texgyretermes-italic.otf]}
    \def\TKRomanBold{[texgyretermes-bold.otf]}
    \def\TKRomanBoldItalic{[texgyretermes-bolditalic.otf]}
  \else
    \def\TKRomanRegular{Times New Roman/OT}
    \def\TKRomanItalic{Times New Roman Italic/OT}
    \def\TKRomanBold{Times New Roman Bold/OT}
    \def\TKRomanBoldItalic{Times New Roman Bold Italic/OT}
  \fi
  \DeclareFontFamily{TU}{ptm}{}
  \DeclareFontShape{TU}{ptm}{m}{n}{<->"\TKRomanRegular:mapping=tex-text"}{}
  \DeclareFontShape{TU}{ptm}{m}{it}{<->"\TKRomanItalic:mapping=tex-text"}{}
  \DeclareFontShape{TU}{ptm}{bx}{n}{<->"\TKRomanBold:mapping=tex-text"}{}
  \DeclareFontShape{TU}{ptm}{bx}{it}{<->"\TKRomanBoldItalic:mapping=tex-text"}{}
  \DeclareFontShape{TU}{ptm}{b}{n}{<->ssub*ptm/bx/n}{}
  \DeclareFontShape{TU}{ptm}{b}{it}{<->ssub*ptm/bx/it}{}
  \DeclareFontShape{TU}{ptm}{m}{sl}{<->ssub*ptm/m/it}{}
  \DeclareFontShape{TU}{ptm}{bx}{sl}{<->ssub*ptm/bx/it}{}
  \ifTKTermes
    \DeclareFontShape{TU}{ptm}{m}{sc}{<->"\TKRomanRegular:mapping=tex-text;+smcp"}{}
  \else
    \DeclareFontShape{TU}{ptm}{m}{sc}{<->ssub*ptm/m/n}{}
  \fi
\fi

\documentclass[journal]{IEEEtran}
\usepackage{cite}
\usepackage{amsmath,amssymb,amsfonts}
% Must follow amsmath: Tectonic's cached XeTeX lacks cmss10.
\ifdefined\XeTeXversion
  \DeclareFontFamily{TU}{qhv}{}
  \IfFileExists{texgyreheros-regular.otf}{
    \DeclareFontShape{TU}{qhv}{m}{n}{<->"[texgyreheros-regular.otf]"}{}
  }{
    \DeclareFontShape{TU}{qhv}{m}{n}{<->"Arial"}{}
  }
  \DeclareMathAlphabet{\mathsf}{TU}{qhv}{m}{n}
\fi

\usepackage{xcolor}
\usepackage[utf8]{inputenc}
\usepackage[T1]{fontenc}
\usepackage{microtype,amsthm,algorithm,algorithmic,float}
\usepackage[hidelinks]{hyperref}
\interdisplaylinepenalty=2500
\emergencystretch=2em
\newcommand{\vphi}{\boldsymbol{\phi}}
\newcommand{\vp}{\mathbf{p}}
\newcommand{\vb}{\mathbf{b}}
\newcommand{\vx}{\mathbf{x}}
\newcommand{\vy}{\mathbf{y}}
\newcommand{\ve}{\mathbf{e}}
\newcommand{\vr}{\mathbf{r}}
\newcommand{\vu}{\mathbf{u}}
\newcommand{\vs}{\mathbf{s}}
\newcommand{\vz}{\mathbf{z}}
\newcommand{\vc}{\mathbf{c}}
\newcommand{\mH}{\mathbf{H}}
\newcommand{\mA}{\mathbf{A}}
\newcommand{\mB}{\mathbf{B}}
\newcommand{\mK}{\mathbf{K}}
\newcommand{\mU}{\mathbf{U}}
\newcommand{\mV}{\mathbf{V}}
\newcommand{\mW}{\mathbf{W}}
\newcommand{\mX}{\mathbf{X}}
\newcommand{\R}{\mathbb{R}}
\newcommand{\norm}[1]{\left\lVert #1\right\rVert}
\newcommand{\col}{\operatorname{col}}
\theoremstyle{plain}
\newtheorem{theorem}{Theorem}
\newtheorem{lemma}{Lemma}
\newtheorem{proposition}{Proposition}
\theoremstyle{definition}
\newtheorem{assumption}{Assumption}
\newtheorem{remark}{Remark}
\theoremstyle{definition}
\newtheorem{definition}{Definition}
\newcommand{\envend}{\leavevmode\unskip\penalty9999 \hbox{}\nobreak\hfill\quad\hbox{\IEEEQEDopen}}

\begin{document}
\title{Transformer--Koopman: History-Conditioned Prediction and Frozen-LTV MPC for Dynamic Trajectory Tracking}
\author{Bolin Tao}
\markboth{IEEE Journal Draft}{Tao: Transformer--Koopman}
\maketitle

\setcounter{section}{1} % Retain the original II--IV section numbering.

\section{Preliminaries and Problem Formulation}
\label{sec:prelim}

\subsection{Koopman Operator and Data-Driven Approximation}
\label{subsec:koopman_prelim}

Consider the discrete-time nonlinear system
\begin{equation}
\vs_{k+1}=f(\vs_k,\vu_k),\qquad
\vx_k=h(\vs_k),
\label{eq:nonlinear_system}
\end{equation}
where $\vs_k\in\mathcal S\subseteq\R^{n_s}$ is the physical state, $\vu_k\in\R^m$ is the control input, and $\vx_k\in\R^n$ is the measured output. The set $\mathcal S$ is the state space of physically admissible states. The map $f:\mathcal S\times\R^m\to\mathcal S$ is the state transition map, and the map $h:\mathcal S\to\R^n$ is the measurement map. In this work, we do not assume direct access to the physical state, and both the predictor and the controller rely on the available input--output data.

The Koopman operator acts on observables, which are scalar-valued
functions of the physical state~\cite{koopman1931hamiltonian,mezic2005spectral,brunton2022modern}. Formally, an observable is defined as
\begin{equation}
\varphi:\mathcal S\to\R.
\label{eq:observable}
\end{equation}
For a given control input, the controlled Koopman operator
$\mathcal K$ maps an observable $\varphi$ to a new function
whose value at the current state $\vs_k$ equals the value
of $\varphi$ at the successor state,
\begin{equation}
(\mathcal K\varphi)(\vs_k)
=\varphi\!\left(f(\vs_k,\vu_k)\right)
=\varphi(\vs_{k+1}).
\label{eq:koopman_operator}
\end{equation}
The operator $\mathcal K$ is linear with respect to the
observables and satisfies
\begin{equation}
\mathcal K(\alpha\varphi+\beta\psi)
=\alpha\mathcal K\varphi+\beta\mathcal K\psi
\label{eq:koopman_linearity}
\end{equation}
for any two observables $\varphi$ and $\psi$ and any
scalars $\alpha$ and $\beta$.

Since the Koopman operator generally acts on an infinite-dimensional
function space, its numerical implementation requires a
finite-dimensional approximation constructed from trajectory data. However, under partial observation, the physical state is not directly available, and the construction of instantaneous observables is therefore restricted to functions of the measured output $\vx=h(\vs)$. Accordingly, we introduce a collection of
scalar-valued functions $\{g_j:\R^n\to\R\}_{j=1}^{d}$ defined
on the measurement space and construct the instantaneous lifting
map as
\begin{equation}
\vz_k=g(\vx_k)
=[g_1(\vx_k),\ldots,g_d(\vx_k)]^\top\in\R^d.
\label{eq:instantaneous_lifting}
\end{equation}
The resulting lifted state provides a finite-dimensional
representation of the system in the selected observable coordinates, without requiring exact reconstruction of the physical state or assuming that the selected observable subspace is invariant under the Koopman operator.

To identify the finite-dimensional Koopman approximation from data,
we collect a set of measured input--output trajectories given by
\begin{equation}
\mathcal D=\{\tau^{(j)}\}_{j=1}^{N_{\rm tr}},\qquad
\tau^{(j)}=\{(\vx_t^{(j)},\vu_t^{(j)})\}_{t=0}^{T_j},
\label{eq:training_data}
\end{equation}
where $N_{\rm tr}$ denotes the number of trajectories and $T_j$
denotes the final time index of the $j$-th trajectory.
From these trajectories, we extract a total of
$M=\sum_{j=1}^{N_{\rm tr}}T_j$ one-step transition samples,
indexed by $(\vx_\ell^-,\vu_\ell,\vx_\ell^+)$ for
$\ell=1,\ldots,M$, where $\vx_\ell^-$ and $\vx_\ell^+$ denote
the measurements before and after the application of the control
input $\vu_\ell$, respectively.
By arranging these samples columnwise, we construct the snapshot
matrices
\begin{equation}
\begin{aligned}
\Gamma   &= [\vx_1^-,\ldots,\vx_M^-],\\
\Gamma^+ &= [\vx_1^+,\ldots,\vx_M^+],\\
\Upsilon &= [\vu_1,\ldots,\vu_M],
\end{aligned}
\label{eq:edmd_snapshot_matrices}
\end{equation}
where $\Gamma,\Gamma^+\in\R^{n\times M}$ and
$\Upsilon\in\R^{m\times M}$.
To represent the system evolution in the lifted space, we apply
the lifting map $g$ to each column of $\Gamma$ and $\Gamma^+$
and obtain the lifted snapshot matrices
\begin{equation}
\begin{aligned}
Z   &= [g(\vx_1^-),\ldots,g(\vx_M^-)],\\
Z^+ &= [g(\vx_1^+),\ldots,g(\vx_M^+)],
\end{aligned}
\label{eq:edmd_lifted_matrices}
\end{equation}
where $Z,Z^+\in\R^{d\times M}$.

Given the lifted snapshot matrices and a prescribed lifting map
$g$, we can identify an input-affine finite-dimensional Koopman
model using extended dynamic mode decomposition with control
(EDMDc). Specifically, the model matrices are obtained by solving
the following least-squares problem
\cite{williams2015data,proctor2016dynamic}:
\begin{equation}
\mK^*=\underset{\mK}{\operatorname{arg\,min}}\;
\left\|Z^+-\mK
\begin{bmatrix}Z\\\Upsilon\end{bmatrix}\right\|_F^2,
\label{eq:edmd_control_optimization}
\end{equation}
where $\mK=[\mA\ \mB]\in\R^{d\times(d+m)}$.
The minimum-Frobenius-norm least-squares solution is given by
\begin{equation}
\mK^*=Z^+\begin{bmatrix}Z\\\Upsilon\end{bmatrix}^{\dagger},
\label{eq:edmd_control_solution}
\end{equation}
where $\dagger$ denotes the Moore--Penrose pseudoinverse.
By partitioning $\mK^*$ into the state transition matrix $\mA$
and the control input matrix $\mB$, the lifted dynamics can be
expressed as
\begin{equation}
\vz_{k+1}=\mA\vz_k+\mB\vu_k+\vr_k.
\label{eq:linear_lifted}
\end{equation}
Here, $\mA\in\R^{d\times d}$ approximates the evolution of the
selected observables, and $\mB\in\R^{d\times m}$ characterizes
the contribution of the control input to the lifted dynamics.
Together, these matrices provide a finite-dimensional
approximation of the controlled Koopman evolution in the
selected observable coordinates, with $\vr_k\in\R^d$
accounting for the one-step prediction error of this
approximation.

\subsection{History-Based Prediction and Problem Formulation}
\label{sec:problem}

The finite-dimensional Koopman model in~\eqref{eq:linear_lifted}
has two major limitations. First, under partial observation, the
instantaneous lifting map $g(\vx_k)$ may not capture sufficient
information for predicting the system evolution, as different
physical states may produce identical measurements but exhibit
different dynamics under the same input. Second, the fixed matrices
$\mA$ and $\mB$ may not adequately approximate the lifted dynamics
across different operating conditions. To address these limitations,
we construct a history-dependent lifted representation and introduce
context-dependent transition and input matrices.

Let $P_h\geq2$ denote the history length. At time $k$, the
measurement history and the corresponding past-input history
are defined as
\begin{align}
\mX_k&=\col(\vx_{k-P_h+1},\ldots,\vx_k)\in\R^{P_hn},
\label{eq:measurement_window}\\
v_k&=\col(\vu_{k-P_h+1},\ldots,\vu_{k-1})
\in\R^{(P_h-1)m},
\label{eq:past_input_history}
\end{align}
where $\col(\cdot)$ denotes vertical stacking.

Based on the measurement history, we redefine the lifted state
$\vz_k$ through a history-dependent lifting map $G_\theta$ as
\begin{equation}
\vz_k=G_\theta(\mX_k)
=\col(\vx_k,\psi_\theta(\mX_k))\in\R^d,
\label{eq:state_preserving_lifting}
\end{equation}
where $d=n+n_\psi$ and $\psi_\theta(\mX_k)\in\R^{n_\psi}$
represents the learned predictive features.
Unlike the instantaneous lifting map $g$, the new lifting map
$G_\theta$ explicitly retains the current measurement and
augments it with features extracted from the measurement history.
This structure allows the measured output to be recovered through
a fixed linear readout,
\begin{equation}
\vx_k=C_x\vz_k,\qquad C_x=[I_n\ 0].
\label{eq:state_readout}
\end{equation}

Given the available histories $(\mX_k,v_k)$ and a prediction
horizon $N\geq1$, this work considers two related problems:
\begin{itemize}
    \item \emph{Multi-step prediction:}
    Given a candidate input sequence $\vu_{k:k+N-1}$,
    predict the future measurements
    $\vx_{k+1},\ldots,\vx_{k+N}$.

    \item \emph{Reference-tracking control:}
    Determine an admissible input sequence such that the system
    outputs track the reference trajectory
    $\vx_{k+1}^r,\ldots,\vx_{k+N}^r$,
    subject to $\vu_{i|k}\in\mathcal U$ for $i=0,\ldots,N-1$,
    where $\mathcal U\subset\R^m$ is a nonempty, compact,
    convex polytope.
\end{itemize}

\section{History-Conditioned Scheduled Koopman Modeling}
\label{sec:model}

\subsection{History-Conditioned Koopman Predictor}
\label{sec:temporal_lifting}

To address the two limitations identified in the previous section,
we develop a history-conditioned scheduled Koopman predictor that
combines history-dependent lifting with context-dependent lifted
dynamics. The former extracts predictive features from the
measurement history, while the latter adapts the lifted matrices
to the operating condition described by the available histories.

To implement the feature map $\psi_\theta$ in the history-dependent
lifting, we employ a patch-based Transformer encoder to capture
temporal dependencies within the measurement history. Specifically,
the measurement window $\mX_k$ is arranged as a $P_h\times n$
array and divided into $N_p=P_h/p$ nonoverlapping temporal patches
of length $p$, where $p$ divides $P_h$. Each patch
$\mX_k^{(j)}\in\R^{p\times n}$ is mapped to a token through
a shared affine embedding,
\begin{equation}
\vphi_k^{(j)}
=\mW_p\operatorname{vec}(\mX_k^{(j)})+\vb_p
\in\R^{d_{\rm model}},
\label{eq:patch_embedding}
\end{equation}
where $j=1,\ldots,N_p$, $\operatorname{vec}(\cdot)$ denotes
vectorization, and $\mW_p\in\R^{d_{\rm model}\times pn}$
and $\vb_p\in\R^{d_{\rm model}}$ are learnable parameters
shared across all patches.

To incorporate temporal order, sinusoidal positional encodings
$\vp^{(j)}\in\R^{d_{\rm model}}$ \cite{vaswani2017attention}
are added to the patch tokens before they are processed by an
$L_T$-layer Transformer encoder,
\begin{equation}
\mH_k=\operatorname{TransformerEncoder}
\left(\{\vphi_k^{(j)}+\vp^{(j)}\}_{j=1}^{N_p}\right),
\label{eq:transformer_encoder}
\end{equation}
where $\mH_k\in\R^{N_p\times d_{\rm model}}$ contains the
contextualized patch representations. A compact representation
of the measurement history is then obtained by mean pooling,
\begin{equation}
\ve_k^x=\frac{1}{N_p}\sum_{j=1}^{N_p}\mH_k^{(j)}
\in\R^{d_{\rm model}},
\label{eq:pooled_history_feature}
\end{equation}
where $\mH_k^{(j)}$ denotes the $j$th patch representation.
Based on the extracted history feature, we implement the
history-dependent lifting map $G_\theta$ as
\begin{equation}
\vz_k^E=G_\theta(\mX_k)
=\col(\vx_k,\mW_h\ve_k^x+\vb_h)\in\R^d,
\label{eq:lifting_head}
\end{equation}
where $\mW_h\in\R^{(d-n)\times d_{\rm model}}$ and
$\vb_h\in\R^{d-n}$ define the linear projection associated
with $\psi_\theta$, and $\theta$ collects all parameters
of the lifting network. The superscript $E$ distinguishes
the encoded lifted state from the recursively propagated
state introduced below.

To overcome the limitations of fixed lifted matrices, we further
introduce a scheduling context constructed from the measurement
and past-input histories. While $\ve_k^x$ characterizes the
measurement history, the past-input history $v_k$ provides
complementary information about the preceding control actions.
We encode $v_k$ using a multilayer perceptron (MLP),
\begin{equation}
\ve_k^u=M_\eta(v_k)\in\R^{d_u},
\label{eq:input_history_feature}
\end{equation}
and combine the two history features through a fusion MLP
$H_\phi$ to obtain the scheduling context,
\begin{equation}
\vc_k=H_\phi\left(\col(\ve_k^x,\ve_k^u)\right)
\in\R^{d_c}.
\label{eq:fused_context}
\end{equation}
The resulting context encoder is denoted by $C$.
It combines the shared measurement encoder of $G_\theta$,
the input encoder $M_\eta$, and the fusion network $H_\phi$.
Since the context is constructed exclusively from the available
histories, it does not explicitly depend on the current control
input $\vu_k$.

Rather than directly generating the full lifted matrices from
the scheduling context, we parameterize them as context-dependent
low-rank corrections to nominal matrices. This parameterization
limits the rank of the matrix variations and enables the model
to adapt its lifted dynamics without generating all matrix
entries independently,
\begin{align}
\mA(\vc)&=\mA_0+
\mU_A\operatorname{diag}(\gamma_A(\vc))\mV_A^\top,
\label{eq:scheduled_A}\\
\mB(\vc)&=\mB_0+
\mU_B\operatorname{diag}(\gamma_B(\vc))\mV_B^\top,
\label{eq:scheduled_B}
\end{align}
where $\mA_0\in\R^{d\times d}$ and $\mB_0\in\R^{d\times m}$
are the nominal lifted matrices. The learnable factors
$\mU_A,\mV_A\in\R^{d\times r_A}$,
$\mU_B\in\R^{d\times r_B}$, and $\mV_B\in\R^{m\times r_B}$
determine the correction directions, while
$\gamma_A(\vc)\in\R^{r_A}$ and $\gamma_B(\vc)\in\R^{r_B}$
determine their context-dependent amplitudes. The correction
ranks satisfy $1\leq r_A\leq d$ and $1\leq r_B\leq\min(d,m)$.
The columns of the low-rank factors are constrained to be orthonormal,
$\mU_\chi^\top\mU_\chi=\mV_\chi^\top\mV_\chi=I_{r_\chi}$
for $\chi\in\{A,B\}$.

To regulate the magnitude of the low-rank corrections,
we parameterize the $j$th scheduling coefficient for each
branch $\chi\in\{A,B\}$ as
\begin{equation}
\begin{aligned}
\gamma_{\chi,j}(\vc)
={}&g_\chi\left[\operatorname{softmax}
(\mW_\chi^\pi\vc+\vb_\chi^\pi)\right]_j\\
&\cdot\tanh\left([\mW_\chi\vc+\vb_\chi]_j\right),
\end{aligned}
\label{eq:scheduling_function}
\end{equation}
where $j=1,\ldots,r_\chi$,
$\mW_\chi^\pi,\mW_\chi\in\R^{r_\chi\times d_c}$ and
$\vb_\chi^\pi,\vb_\chi\in\R^{r_\chi}$ are learnable parameters,
and $g_\chi\geq0$ is a scaling gain.
The softmax function distributes the correction weights
across the rank-one components, while the hyperbolic tangent
determines their signed amplitudes. Under the imposed orthonormality constraints, the resulting
correction satisfies
\begin{equation}
\left\|\mU_\chi\operatorname{diag}(\gamma_\chi(\vc))
\mV_\chi^\top\right\|_2
=\max_j|\gamma_{\chi,j}(\vc)|\leq g_\chi.
\label{eq:gain_bound}
\end{equation}
Thus, the gain $g_\chi$ bounds the deviation of the scheduled
matrix from its nominal value, but does not by itself
guarantee the stability of the lifted dynamics.

By integrating the history-dependent lifting map and the
context-dependent lifted matrices, we formulate the proposed
predictor as follows.

\begin{definition}[History-Conditioned Scheduled Koopman Predictor]
\label{def:history_scheduled_predictor}
Given the measured histories $(\mX_k,v_k)$,
the predictor is initialized as
\begin{equation}
\vz_k^E=G_\theta(\mX_k),\qquad
\widehat{\mX}_k=\mX_k.
\label{eq:encoded_initialization}
\end{equation}
Starting from $\vz_k=\vz_k^E$, the scheduling context and
lifted dynamics are given by
\begin{align}
\vc_k&=C(\widehat{\mX}_k,v_k),
\label{eq:context_at_step}\\
\vz_{k+1}&=\mA(\vc_k)\vz_k+\mB(\vc_k)\vu_k.
\label{eq:lifted_recursive_dynamics}
\end{align}
The measurement and input histories are recursively updated as
\begin{equation}
\left\{
\begin{aligned}
\widehat{\mX}_{k+1}
&=S_X\widehat{\mX}_k+E_XC_x\vz_{k+1},\\
v_{k+1}
&=S_uv_k+E_u\vu_k,
\end{aligned}
\right.
\label{eq:history_window_update}
\end{equation}
where $S_X\in\R^{P_hn\times P_hn}$ and
$S_u\in\R^{(P_h-1)m\times(P_h-1)m}$ are shift matrices
that remove the oldest entries from their respective histories,
and $E_X\in\R^{P_hn\times n}$ and
$E_u\in\R^{(P_h-1)m\times m}$ append the latest predicted
measurement and control input, respectively.\envend
\end{definition}

During multi-step prediction, the lifted state is initialized
once through $G_\theta$ and subsequently propagated through
the scheduled lifted dynamics without re-encoding.
In contrast, the scheduling context is recomputed at each
prediction step using the updated measurement and input
histories. Although this requires reevaluating the shared
measurement encoder, the resulting features are used only
to update the scheduling coefficients and do not replace
the propagated lifted state.
Given a candidate input sequence, recursively evaluating
the predictor yields a multi-step lifted trajectory, from
which the future measurements are obtained through the
fixed linear readout $C_x$.

\subsection{Model Learning}
\label{sec:training}

To separate the learning of the lifted representation from
the adaptation of the lifted dynamics, we adopt a two-stage
training strategy using the trajectories in~\eqref{eq:training_data}.
The first stage jointly learns the lifting map and nominal
lifted matrices, while the second stage learns context-dependent
corrections through recursive multi-step prediction within
the established lifted coordinates.

In the first stage, the scheduling corrections are disabled,
and the predictor reduces to a fixed-coefficient Koopman model,
\begin{equation}
\vz_{k+1}=\mA_0\vz_k+\mB_0\vu_k.
\label{eq:nominal_lifted_dynamics}
\end{equation}
Unlike the conventional EDMDc formulation with a prescribed
lifting map, we jointly optimize the history-dependent lifting
map $G_\theta$ and the nominal matrices $\mA_0$ and $\mB_0$.
Given the encoded lifted state $\vz_k^E=G_\theta(\mX_k)$,
the nominal one-step prediction is given by
$\widetilde\vz_{k+1}^{0}=\mA_0\vz_k^E+\mB_0\vu_k$,
where the tilde denotes a predicted lifted state and the
superscript $0$ indicates that the prediction is generated
by the nominal model without scheduling corrections.
The training objective is defined as
\begin{equation}
\begin{split}
\mathcal L_{\rm init}=\sum_{k\in\mathcal I_1}\bigl[&
\norm{C_x\widetilde\vz_{k+1}^0-\vx_{k+1}}_2^2\\
&+\alpha_0\norm{\widetilde\vz_{k+1}^0-\vz_{k+1}^E}_2^2
\bigr],
\end{split}
\label{eq:training_one_step_objects}
\end{equation}
where $\mathcal I_1$ indexes the one-step training samples
$(\mX_k,\vu_k,\mX_{k+1})$, $\alpha_0>0$ is a weighting
coefficient, and $\vz_{k+1}^E=G_\theta(\mX_{k+1})$.
The first term minimizes the prediction error in the measurement
space, while the second promotes consistency between the
predicted and encoded lifted states.

In the second stage, we freeze the lifting map $G_\theta$
and the nominal matrices $\mA_0$ and $\mB_0$, and optimize
only the scheduling parameters, including those of $M_\eta$,
$H_\phi$, the scheduling coefficient networks, and the low-rank
factors $\mU_\chi$ and $\mV_\chi$.
The low-rank factors are optimized subject to the orthonormality
constraints imposed in the preceding subsection.
This strategy preserves the learned lifted coordinates and
allows the scheduling mechanism to focus on adapting the
nominal dynamics to different operating conditions.

To improve multi-step prediction accuracy, the scheduling
mechanism is trained through recursive rollouts over a horizon
$H_R$. Given the measured histories $(\mX_k,v_k)$, each
rollout is initialized with $\vz_{0|k}^P=G_\theta(\mX_k)$,
$\widehat{\mX}_{0|k}^P=\mX_k$, and $v_{0|k}^P=v_k$, where
the superscript $P$ denotes propagated quantities and $i|k$
denotes prediction step $i$ from time $k$.
Under the recorded input sequence $\vu_{k:k+H_R-1}$,
the predictor in Definition~\ref{def:history_scheduled_predictor}
is recursively evaluated to obtain the lifted trajectory
$\{\vz_{i|k}^P\}_{i=1}^{H_R}$.
Extending the one-step objective to multi-step prediction,
we define the rollout loss as
\begin{equation}
\begin{split}
\mathcal L_R=\sum_{k\in\mathcal I_R}\sum_{i=1}^{H_R}
\omega_i\bigl[&
\norm{C_x\vz_{i|k}^P-\vx_{k+i}}_2^2\\
&+\alpha_R\norm{\vz_{i|k}^P-\vz_{k+i}^E}_2^2\bigr],
\end{split}
\label{eq:joint_training_loss}
\end{equation}
where $\mathcal I_R$ indexes the multi-step training samples,
$\omega_i=1/H_R$, and $\alpha_R>0$ balances the measurement
prediction error and lifted-state consistency.
The consistency targets $\vz_{k+i}^E=G_\theta(\mX_{k+i})$
are obtained from the recorded measurement windows, whereas
the scheduling contexts are evaluated using the recursively
updated histories.
In this way, the scheduling mechanism is trained under the
prediction errors encountered during recursive rollouts.
We set $H_R=N$ to align the training horizon with the
prediction horizon of the MPC developed in
Section~\ref{sec:control}.

\subsection{Prediction Error Analysis}
\label{sec:prediction_error_analysis}

To quantify the predictive accuracy of the learned model, we compare
the recursively propagated lifted state with the encoded state along
a trajectory of the physical system. Fix a sampling instant $k$ and a
feasible candidate sequence $U_k=\col(\vu_{0|k},\ldots,\vu_{N-1|k})
\in\mathcal U^N$. In the following analysis, the dependence on $U_k$
is suppressed when unambiguous, and the abbreviated index $i$ refers
to prediction step $i|k$. All vector norms are Euclidean and all matrix
norms are induced $2$-norms unless otherwise specified.

Let $\mX_{i|k}^E$ be the measurement history generated by the physical
system under $U_k$. The corresponding encoded state and scheduling
context are $\vz_{i|k}^E=G_\theta(\mX_{i|k}^E)$ and
$\vc_{i|k}^E=C(\mX_{i|k}^E,v_i)$, respectively. Here the superscript
$E$ denotes quantities encoded from the physical trajectory. Without
assuming exact closure in the learned coordinates, we write
\begin{equation}
\vz_{i+1|k}^E=\mA(\vc_{i|k}^E)\vz_{i|k}^E
+\mB(\vc_{i|k}^E)\vu_i+\vr_{i|k},
\label{eq:encoded_residual_identity}
\end{equation}
where $\vu_i=\vu_{i|k}$ and $\vr_{i|k}$ is the one-step lifted
closure residual defined by this identity. The propagated trajectory
$\vz_{i|k}^P$ starts from $\vz_{0|k}^P=\vz_{0|k}^E=G_\theta(\mX_k)$
and follows Definition~\ref{def:history_scheduled_predictor}.
For compact notation, write $\mX_{i|k}^P=\widehat{\mX}_{i|k}^P$ for the
propagated measurement history and $\vc_{i|k}^P=C(\mX_{i|k}^P,v_i)$.
The two trajectories share the initial measurement and input histories.
Since the past-input history is updated by $v_{i+1}=S_uv_i+E_u\vu_i$
with $v_0=v_k$, it is identical along both trajectories. Their contexts
therefore differ only through the measurement histories.

We first specify the operating domain, the closure residual bound, and
the regularity of the context encoder. The assumptions are understood
uniformly over the initial conditions and reference times in the
operating region of interest.

\begin{assumption}[Bounded history domain]
\label{asm:domain}
There exists a compact set $\mathcal H\subseteq
\R^{P_hn}\times\R^{(P_h-1)m}$ containing the history pairs along the
physical and recursively predicted trajectories for every considered
initial condition, sampling instant $k$, feasible $U_k$, and prediction
step $i=0,\ldots,N$. This includes the nominal rollout introduced in
Section~\ref{sec:frozen_mpc_formulation}.
\end{assumption}

\begin{assumption}[Bounded lifted states and inputs]
\label{asm:bounds}
On $\mathcal H$, $\norm{G_\theta(\mX)}_2\le M_z$. The propagated
lifted states, including those along the nominal rollout, satisfy
$\norm{\vz_{i|k}^P}_2\le M_P$ for some finite $M_P$. The input set $\mathcal U$ is as
specified in Section~\ref{sec:problem}, and
$M_u=\max_{\vu\in\mathcal U}\norm{\vu}_2$.
\end{assumption}

\begin{assumption}[Uniform closure residual]
\label{asm:residual}
There exists $\varepsilon\ge0$ such that
$\norm{\vr_{i|k}}_2\le\varepsilon$ for every physical trajectory
considered in Assumption~\ref{asm:domain} and $i=0,\ldots,N-1$.
\end{assumption}

\begin{assumption}[Lipschitz context encoder]
\label{asm:context_lipschitz}
There exist constants $L_{Cx},L_{Cu}\ge0$ such that, for any
$(\mX,v),(\mX',v')\in\mathcal H$,
\begin{equation}
\begin{aligned}
&\norm{C(\mX,v)-C(\mX',v')}_2\\
&\qquad\le L_{Cx}\norm{\mX-\mX'}_2+L_{Cu}\norm{v-v'}_2.
\end{aligned}
\label{eq:context_lipschitz}
\end{equation}
\end{assumption}

The constants in Assumption~\ref{asm:context_lipschitz} describe the
realized encoder on $\mathcal H$ and depend on its learned parameters
and constituent nonlinear maps. The joint inequality is stated
without requiring $\mathcal H$ to be a Cartesian product. The residual
bound in Assumption~\ref{asm:residual} is an a posteriori modeling
condition rather than a consequence of the training loss. Validation
trajectories can be used to estimate the residual, but a finite-sample
maximum alone does not certify the uniform bound.

The low-rank correction structure provides explicit matrix bounds
and Lipschitz constants, as established in the following lemma.

\begin{lemma}[Regularity of the scheduled matrices]
\label{lem:scheduling_regularity}
Let $\gamma_\chi$ be defined by~\eqref{eq:scheduling_function} for
$\chi\in\{A,B\}$, with the low-rank factors satisfying the
orthonormality constraints of the model. Define
\begin{equation}
\begin{aligned}
a&=\norm{\mA_0}_2+g_A,\\
b_B&=\norm{\mB_0}_2+g_B,\\
L_\chi&=g_\chi\Big(\tfrac12\norm{\mW_\chi^\pi}_2+\norm{\mW_\chi}_2\Big).
\end{aligned}
\label{eq:regularity_constants}
\end{equation}
Then the norm identity~\eqref{eq:gain_bound} holds, and for all
$\vc,\vc'\in\R^{d_c}$,
\begin{gather}
\norm{\mA(\vc)}_2\le a,\qquad\norm{\mB(\vc)}_2\le b_B,
\label{eq:uniform_matrix_bound}\\
\begin{aligned}
\norm{\mA(\vc)-\mA(\vc')}_2&\le L_A\norm{\vc-\vc'}_2,\\
\norm{\mB(\vc)-\mB(\vc')}_2&\le L_B\norm{\vc-\vc'}_2.
\end{aligned}
\label{eq:scheduled_lipschitz}
\end{gather}
\end{lemma}

\begin{proof}
The argument proceeds from the correction term to the full matrices and
then to the context dependence. Fix one branch and suppress the
subscript $\chi$; let $\mV_j$ denote the $j$th column of $\mV$. For any
compatible vector $\mathbf w$, the orthonormal columns of $\mU$ imply
$\norm{\mU\mathbf y}_2=\norm{\mathbf y}_2$ for every
$\mathbf y\in\R^{r_\chi}$, so that
\begin{equation*}
\begin{aligned}
\norm{\mU\operatorname{diag}(\gamma)\mV^\top\mathbf w}_2^2
&=\norm{\operatorname{diag}(\gamma)\mV^\top\mathbf w}_2^2
=\sum_{j=1}^{r_\chi}\gamma_j^2(\mV_j^\top\mathbf w)^2\\
&\le\max_j\gamma_j^2\,\norm{\mV^\top\mathbf w}_2^2
\le\max_j\gamma_j^2\,\norm{\mathbf w}_2^2,
\end{aligned}
\end{equation*}
where the last step uses $\norm{\mV^\top}_2=1$, which follows from the
orthonormal columns of $\mV$.
Equality is attained by $\mathbf w=\mV_j$ at an index maximizing
$|\gamma_j|$. Since every softmax entry lies in $[0,1]$ and
$|\tanh(\cdot)|\le1$, each coefficient obeys
$|\gamma_j(\vc)|\le g_\chi$, which proves
\eqref{eq:gain_bound}. Applying the triangle inequality
to~\eqref{eq:scheduled_A}--\eqref{eq:scheduled_B}, with this identity
used for both correction terms, then
yields~\eqref{eq:uniform_matrix_bound}.

It remains to quantify the context dependence. For
$\vp=\operatorname{softmax}(\boldsymbol\zeta)$ with
$\boldsymbol\zeta\in\R^{r_\chi}$, the Jacobian is
$\mathbf J=\operatorname{diag}(\vp)-\vp\vp^\top$, which is symmetric
and positive semidefinite. Its spectral norm therefore equals the
spectral radius, which is bounded by every induced matrix norm. Using
the induced $\infty$-norm, that is, the maximum absolute row sum, and
$\sum_jp_j=1$,
\begin{equation}
\norm{\mathbf J}_2\le\norm{\mathbf J}_\infty
=\max_j2p_j(1-p_j)\le\tfrac12.
\label{eq:softmax_jacobian_bound}
\end{equation}
This bound transfers to the scheduling coefficients as follows. Write
$\gamma(\vc)=g_\chi\vp(\vc)\odot\mathbf t(\vc)$, where $\vp(\vc)$ is
the softmax branch and $\mathbf t(\vc)$ the tanh branch
of~\eqref{eq:scheduling_function}, and $\odot$ denotes elementwise
multiplication. Since both branches are affine in the context before
the nonlinearity, the mean value inequality
with~\eqref{eq:softmax_jacobian_bound} and the unit Lipschitz constant
of $\tanh$ give
\begin{equation*}
\begin{aligned}
\norm{\vp(\vc)-\vp(\vc')}_2
&\le\tfrac12\norm{\mW_\chi^\pi}_2\norm{\vc-\vc'}_2,\\
\norm{\mathbf t(\vc)-\mathbf t(\vc')}_2
&\le\norm{\mW_\chi}_2\norm{\vc-\vc'}_2.
\end{aligned}
\end{equation*}
The product is then split as
$\vp\odot\mathbf t-\vp'\odot\mathbf t'
=(\vp-\vp')\odot\mathbf t+\vp'\odot(\mathbf t-\mathbf t')$, with
primes denoting evaluation at $\vc'$. Because
$\norm{\vp'}_\infty,\norm{\mathbf t}_\infty\le1$, the two estimates
combine into
\begin{equation*}
\begin{aligned}
\norm{\gamma(\vc)-\gamma(\vc')}_2
&\le g_\chi\big(\norm{\vp(\vc)-\vp(\vc')}_2
+\norm{\mathbf t(\vc)-\mathbf t(\vc')}_2\big)\\
&\le L_\chi\norm{\vc-\vc'}_2,
\end{aligned}
\end{equation*}
with $L_\chi$ as in~\eqref{eq:regularity_constants}. Finally, the norm
identity derived above holds for an arbitrary coefficient vector.
Applying it to the coefficient difference
$\Delta\gamma=\gamma(\vc)-\gamma(\vc')$ gives
\begin{equation*}
\norm{\mU\operatorname{diag}(\Delta\gamma)\mV^\top}_2
=\max_j|\Delta\gamma_j|
\le\norm{\Delta\gamma}_2
\le L_\chi\norm{\vc-\vc'}_2,
\end{equation*}
which is~\eqref{eq:scheduled_lipschitz}.
\end{proof}

Lemma~\ref{lem:scheduling_regularity} shows that the correction gains
govern both the magnitude and the context sensitivity of the scheduled
dynamics. The statement concerns norms only and does not imply Schur
stability of the lifted dynamics. It does, however, supply a
finite-horizon bound on the propagated state, namely the choice
$M_P=\max_{0\le i\le N}(a^iM_z+b_BM_u\sum_{j=0}^{i-1}a^j)$ in
Assumption~\ref{asm:bounds}, with the empty sum equal to zero.
To quantify the recursive prediction error, define
\begin{equation}
b=L_{Cx}(L_AM_z+L_BM_u),\qquad
\rho=a+b\sqrt{P_h}.
\label{eq:rho_definition}
\end{equation}

\begin{theorem}[Multi-step prediction error bound]
\label{thm:prediction_error}
Under Assumptions~\ref{asm:domain}--\ref{asm:context_lipschitz}, for
any feasible $U_k$, with $\vz_{0|k}^E=\vz_{0|k}^P$ and coincident input
histories as specified above, and $i=0,\ldots,N$,
\begin{equation}
\norm{\vx_{i|k}^E-\vx_{i|k}^P}_2
\le\norm{\vz_{i|k}^E-\vz_{i|k}^P}_2
\le\mathcal E_i:=\varepsilon\sum_{j=0}^{i-1}\rho^j,
\label{eq:prediction_error_bound}
\end{equation}
where $\vx_{i|k}^E=C_x\vz_{i|k}^E$,
$\vx_{i|k}^P=C_x\vz_{i|k}^P$, and the empty sum gives $\mathcal E_0=0$.
Equivalently,
\begin{equation}
\mathcal E_i=\begin{cases}
\varepsilon(1-\rho^i)/(1-\rho),&\rho\ne1,\\
i\varepsilon,&\rho=1.
\end{cases}
\label{eq:prediction_error_explicit}
\end{equation}
\end{theorem}

\begin{proof}
The proof derives a recursion for the lifted error, bounds the
context mismatch through the measurement-history window, and closes
the resulting inequality with a scalar envelope. Let
$\ve_i=\vz_{i|k}^E-\vz_{i|k}^P$. Subtracting the propagated
dynamics from~\eqref{eq:encoded_residual_identity} gives
\begin{equation}
\begin{aligned}
\ve_{i+1}={}&\mA(\vc_i^P)\ve_i
+\big[\mA(\vc_i^E)-\mA(\vc_i^P)\big]\vz_i^E\\
&+\big[\mB(\vc_i^E)-\mB(\vc_i^P)\big]\vu_i+\vr_i.
\end{aligned}
\label{eq:error_recursion}
\end{equation}
The first term propagates the existing error, the two bracketed terms
arise from the context mismatch, and $\vr_i$ is the closure residual.
The context mismatch is driven by the measurement histories alone:
since the input histories coincide,
Assumption~\ref{asm:context_lipschitz} yields
$\norm{\vc_i^E-\vc_i^P}_2\le L_{Cx}\norm{\mX_i^E-\mX_i^P}_2$.
The initial measurement histories coincide as well, so window entries
with index $i-j<0$ refer to samples taken before $k$ and are identical
along both trajectories. Adopting the convention $\ve_j=0$ for $j<0$
and using the coordinate selection matrix $C_x$, with
$\norm{C_x}_2=1$, we obtain
\begin{equation}
\norm{\mX_i^E-\mX_i^P}_2^2
=\sum_{j=0}^{P_h-1}\norm{C_x\ve_{i-j}}_2^2
\le\sum_{j=0}^{P_h-1}\norm{\ve_{i-j}}_2^2.
\label{eq:history_error_bound}
\end{equation}
We now bound the four terms of~\eqref{eq:error_recursion} in turn.
By Lemma~\ref{lem:scheduling_regularity}, the propagated term satisfies
$\norm{\mA(\vc_i^P)\ve_i}_2\le a\norm{\ve_i}_2$, while
Assumption~\ref{asm:residual} gives $\norm{\vr_i}_2\le\varepsilon$.
For the two context-dependent terms, the Lipschitz bounds
\eqref{eq:scheduled_lipschitz} combine with
$\norm{\vz_i^E}_2\le M_z$ and $\norm{\vu_i}_2\le M_u$ from
Assumption~\ref{asm:bounds} to give
\begin{equation*}
\begin{aligned}
\norm{\big[\mA(\vc_i^E)-\mA(\vc_i^P)\big]\vz_i^E}_2
&\le L_AM_z\norm{\vc_i^E-\vc_i^P}_2,\\
\norm{\big[\mB(\vc_i^E)-\mB(\vc_i^P)\big]\vu_i}_2
&\le L_BM_u\norm{\vc_i^E-\vc_i^P}_2.
\end{aligned}
\end{equation*}
Inserting the context estimate shows that their sum is at most
$b\norm{\mX_i^E-\mX_i^P}_2$, with $b$ defined
in~\eqref{eq:rho_definition}. Together
with~\eqref{eq:history_error_bound}, this yields
\begin{equation}
\norm{\ve_{i+1}}_2\le a\norm{\ve_i}_2+\varepsilon
+b\Big(\sum_{j=0}^{P_h-1}\norm{\ve_{i-j}}_2^2\Big)^{1/2}.
\label{eq:error_inequality}
\end{equation}

The inequality couples the current error with the whole history
window. To close it, define the envelope $\mathcal E_i$ by $\mathcal E_i=0$ for $i\le0$
and $\mathcal E_{i+1}=\rho \mathcal E_i+\varepsilon$ for $i\ge0$. The envelope is
nondecreasing, because $\mathcal E_{i+1}-\mathcal E_i=\varepsilon\rho^i\ge0$. The shared
initialization gives $\ve_0=0$. Suppose that $\norm{\ve_j}_2\le \mathcal E_j$
for all $j\le i$. Every window entry is then bounded by $\mathcal E_i$, and
substituting $\mathcal E_{i-j}\le \mathcal E_i$ into~\eqref{eq:error_inequality} yields
$\norm{\ve_{i+1}}_2\le(a+b\sqrt{P_h})\mathcal E_i+\varepsilon=\mathcal E_{i+1}$.
The lifted error bound follows by induction, and solving the scalar
recurrence gives~\eqref{eq:prediction_error_explicit}. Finally,
$\norm{C_x\ve_i}_2\le\norm{\ve_i}_2$ transfers the bound to the
measurement error.
\end{proof}

Theorem~\ref{thm:prediction_error} separates the one-step closure
residual from its amplification through the nominal dynamics,
context-dependent matrices, and history updates. No contraction is
required over a finite horizon. If $\rho<1$ and the domain and residual
assumptions hold over arbitrarily long rollouts, the bound further
reduces to $\mathcal E_i\le\varepsilon/(1-\rho)$ uniformly in $i$. The bound is
a conditional prediction guarantee and does not itself enforce the
operating-domain assumptions.

\begingroup
\section{Frozen-LTV MPC for Dynamic Trajectory Tracking}
\label{sec:control}

\subsection{Formulation of the Frozen-LTV MPC}
\label{sec:frozen_mpc_formulation}

The history-conditioned predictor developed in Section~\ref{sec:model}
is affine in the current control input for a given scheduling
context. However, during multi-step prediction, the scheduling
context evolves with the propagated measurement and input
histories and consequently depends on the candidate control
sequence. Directly incorporating the recursive predictor into
MPC would therefore generally result in a nonlinear
optimization problem.

To retain the computational structure of quadratic programming,
we adopt a single-freeze strategy in which the scheduling
matrices are evaluated along a feasible nominal trajectory
and subsequently held fixed throughout the MPC optimization.
This strategy allows the lifted dynamics to vary across the
prediction horizon while preserving a linear prediction model
with respect to the optimization variables.

At sampling instant $k$, given the available histories
$(\mX_k,v_k)$ and the prediction horizon $N$, we introduce
a feasible nominal input sequence
\begin{equation}
\bar U_k=\col(\bar\vu_{0|k},\ldots,\bar\vu_{N-1|k})
\in\mathcal U^N,
\label{eq:nominal_input_sequence}
\end{equation}
where the overbar denotes nominal quantities. The initial nominal
sequence $\bar U_0$ can be any feasible sequence in
$\mathcal U^N$, and the subsequent sequences are constructed
from the preceding MPC solutions.

Based on the available histories, the nominal rollout
is initialized as
\begin{equation}
\bar\vz_{0|k}=G_\theta(\mX_k),\quad
\bar\mX_{0|k}=\mX_k,\quad
\bar v_{0|k}=v_k.
\label{eq:nominal_initialization}
\end{equation}

Following the recursive predictor in
Definition~\ref{def:history_scheduled_predictor},
the nominal scheduling context at each prediction step
$i=0,\ldots,N-1$ is computed as
\begin{equation}
\bar\vc_{i|k}=C(\bar\mX_{i|k},\bar v_{i|k}),
\label{eq:nominal_context}
\end{equation}
which determines the corresponding scheduled matrices,
\begin{equation}
\bar\mA_{i|k}=\mA(\bar\vc_{i|k}),\quad
\bar\mB_{i|k}=\mB(\bar\vc_{i|k}).
\label{eq:nominal_matrices}
\end{equation}
The nominal lifted state and history windows are then
recursively propagated according to
\begin{equation}
\left\{
\begin{aligned}
\bar\vz_{i+1|k}
&=\bar\mA_{i|k}\bar\vz_{i|k}
+\bar\mB_{i|k}\bar\vu_{i|k},\\
\bar\mX_{i+1|k}
&=S_X\bar\mX_{i|k}+E_XC_x\bar\vz_{i+1|k},\\
\bar v_{i+1|k}
&=S_u\bar v_{i|k}+E_u\bar\vu_{i|k}.
\end{aligned}
\right.
\label{eq:nominal_rollout}
\end{equation}
In this way, the scheduling matrices are evaluated along
the entire nominal trajectory rather than at the current
sampling instant alone.

After completing the nominal rollout, we freeze the resulting
matrix sequence
$\{\bar\mA_{i|k},\bar\mB_{i|k}\}_{i=0}^{N-1}$
and construct a linear time-varying (LTV) prediction model
in the lifted coordinates,
\begin{equation}
\vz_{0|k}=G_\theta(\mX_k),\qquad
\vz_{i+1|k}=\bar\mA_{i|k}\vz_{i|k}+\bar\mB_{i|k}\vu_{i|k},
\label{eq:frozen_model}
\end{equation}
for $i=0,\ldots,N-1$. Unlike the nominal inputs used to
determine the scheduling matrices, the inputs $\vu_{i|k}$
are decision variables to be optimized. Since the scheduling
matrices are fixed during optimization, neither the history
windows nor the scheduling contexts need to be updated
within the optimization problem.

Let $\vx_{k+i}^r$, $i=1,\ldots,N$, denote the desired
reference measurements, where the superscript $r$ denotes
reference quantities. We define the stacked control sequence as
\begin{equation}
U_k=\col(\vu_{0|k},\ldots,\vu_{N-1|k})\in\R^{Nm}.
\label{eq:stacked_control_sequence}
\end{equation}
Using the fixed linear readout $C_x$, we formulate the
finite-horizon tracking objective directly in the measurement space,
\begin{equation}
J_k(U_k)=\sum_{i=1}^{N}
\norm{C_x\vz_{i|k}-\vx_{k+i}^r}_{Q_x}^2
+\sum_{i=0}^{N-1}\norm{\vu_{i|k}}_{R_u}^2,
\label{eq:quadratic_template}
\end{equation}
where $Q_x\in\R^{n\times n}$ and $R_u\in\R^{m\times m}$
are positive definite weighting matrices, and
$\norm{a}_W^2=a^\top Wa$ denotes the squared weighted norm.
The first term penalizes measurement tracking errors, and the
second penalizes the magnitude of the control inputs.

The resulting frozen-LTV MPC problem is formulated as
\begin{equation}
\begin{aligned}
U_k^*
={}&\underset{U_k}{\operatorname{arg\,min}}\;J_k(U_k)\\
\text{s.t.}\quad
&\vz_{0|k}=G_\theta(\mX_k),\\
&\vz_{i+1|k}
=\bar\mA_{i|k}\vz_{i|k}+\bar\mB_{i|k}\vu_{i|k},\\
&\vu_{i|k}\in\mathcal U,\quad i=0,\ldots,N-1.
\end{aligned}
\label{eq:mpc_problem}
\end{equation}

With the scheduling matrices fixed, the predicted lifted
states are affine functions of $U_k$. Consequently, the
objective is quadratic and the admissible input sequences
form the compact convex polytope $\mathcal U^N$.
Since $Q_x\succ0$ and $R_u\succ0$, the objective is strictly
convex with respect to $U_k$. Together with the nonemptiness
and compactness of $\mathcal U^N$, this ensures that the
resulting QP admits a unique optimal control sequence.

Here, the superscript $*$ denotes optimal quantities.
Following the receding-horizon principle, only the first
optimal input is applied to the system,
\begin{equation}
\vu_k=\vu_{0|k}^*.
\label{eq:apply_first}
\end{equation}
The nominal input sequence for the next sampling instant
is constructed by shifting the optimal sequence forward
and repeating its terminal input,
\begin{equation}
\bar U_{k+1}
=\col(\vu_{1|k}^*,\ldots,\vu_{N-1|k}^*,\vu_{N-1|k}^*).
\label{eq:nominal_shift}
\end{equation}
For $N=1$, we simply set $\bar U_{k+1}=\vu_{0|k}^*$.
Since every optimal input satisfies the input constraints,
the updated nominal sequence remains feasible.

At the next sampling instant, the predictor is reinitialized
using the updated measured histories, and a new sequence
of scheduling matrices is generated. The complete procedure
is summarized in Algorithm~\ref{alg:mpc}.

\begin{algorithm}[H]
\caption{Single-Freeze Transformer--Koopman LTV-MPC}
\label{alg:mpc}
\begin{algorithmic}[1]
\STATE \textbf{Given:} $G_\theta$, $C$, $\mA(\cdot)$,
$\mB(\cdot)$, $\mathcal U$, $N$, $Q_x$, and $R_u$.
\STATE \textbf{Initialize:} $\bar U_0\in\mathcal U^N$.
\FOR{$k=0,1,\ldots$}
\STATE Construct the measured histories $(\mX_k,v_k)$.
\STATE Initialize the nominal rollout using
\eqref{eq:nominal_initialization}.
\STATE Perform the nominal rollout using
\eqref{eq:nominal_context}--\eqref{eq:nominal_rollout}.
\STATE Freeze
$\{\bar\mA_{i|k},\bar\mB_{i|k}\}_{i=0}^{N-1}$.
\STATE Obtain the reference measurements
$\vx_{k+1}^r,\ldots,\vx_{k+N}^r$.
\STATE Solve the quadratic program~\eqref{eq:mpc_problem}.
\STATE Apply the first optimal input using
\eqref{eq:apply_first}.
\STATE Update the nominal input sequence using
\eqref{eq:nominal_shift}.
\ENDFOR
\end{algorithmic}
\end{algorithm}

The single-freeze strategy captures variations in the lifted
dynamics through one nominal rollout per sampling instant while
retaining a convex QP. Section~\ref{sec:freezing_error} quantifies the
additional prediction error due to freezing, and
Section~\ref{sec:closed_loop_stability} uses it to establish
sufficient conditions for practical closed-loop output tracking.

\subsection{Prediction Error Induced by Scheduling Freezing}
\label{sec:freezing_error}

Freezing the scheduling matrices introduces an additional prediction
error beyond the closure residual analyzed in
Section~\ref{sec:prediction_error_analysis}. To isolate this error,
let $\vz_{i|k}^F$ denote the frozen prediction in~\eqref{eq:frozen_model}
under a candidate sequence $U_k$, where the superscript $F$ denotes
frozen quantities. In contrast, $\vz_{i|k}^P$ is propagated with the
context recomputed at each step under the same $U_k$, while
$\bar\vz_{i|k}$ is propagated under $\bar U_k$. All three trajectories
start from the same encoded state and histories.

The frozen matrices are evaluated along the nominal rollout. The
freezing error is therefore governed by the distance between the
recursive trajectory under $U_k$ and the nominal one, which involves
the lifted state and both history windows. To measure this distance,
we collect these quantities in the augmented state
$q=\col(\vz,\mX,v)$ and define the lifted update
$f_z(q,\vu)=\mA(C(\mX,v))\vz+\mB(C(\mX,v))\vu$. The associated
augmented transition map is
\begin{equation}
\mathcal F(q,\vu)=\begin{bmatrix}
f_z(q,\vu)\\ S_X\mX+E_XC_xf_z(q,\vu)\\ S_uv+E_u\vu
\end{bmatrix}.
\label{eq:augmented_update_map}
\end{equation}
Here $q$ belongs to the predictor's bounded augmented domain
$\mathcal Q=\{\col(\vz,\mX,v):(\mX,v)\in\mathcal H,
\norm{\vz}_2\le M_P\}$. We next show that $\mathcal F$ is Lipschitz on
$\mathcal Q$. Applying Cauchy--Schwarz
to~\eqref{eq:context_lipschitz} bounds the context difference by
$L_{Cq}\norm{\col(\mX-\mX',v-v')}_2$, where
\begin{equation}
\begin{aligned}
L_{Cq}&=\sqrt{L_{Cx}^2+L_{Cu}^2},\\
\beta_C&=(L_AM_P+L_BM_u)L_{Cq}.
\end{aligned}
\label{eq:augmented_sensitivity}
\end{equation}
The constant $\beta_C$ measures the sensitivity of the lifted update
to the histories through the scheduled matrices. Indeed, for any
$q,q'\in\mathcal Q$ and $\vu\in\mathcal U$,
Lemma~\ref{lem:scheduling_regularity} gives
\begin{equation*}
\begin{aligned}
\norm{f_z(q,\vu)-f_z(q',\vu)}_2
&\le a\norm{\vz-\vz'}_2\\
&\quad+\beta_C\norm{\col(\mX-\mX',v-v')}_2\\
&\le\sqrt{a^2+\beta_C^2}\norm{q-q'}_2.
\end{aligned}
\end{equation*}
For fixed $q$, the update $f_z$ is affine in $\vu$, so that
$\norm{f_z(q',\vu)-f_z(q',\vu')}_2\le b_B\norm{\vu-\vu'}_2$. To
transfer these two estimates to $\mathcal F$, write
$\Delta_qf=f_z(q,\vu)-f_z(q',\vu)$ and
$\Delta_uf=f_z(q',\vu)-f_z(q',\vu')$, and split the difference
of~\eqref{eq:augmented_update_map} into three stacked vectors,
\begin{equation*}
\begin{aligned}
&\mathcal F(q,\vu)-\mathcal F(q',\vu')
=\begin{bmatrix}0\\ S_X(\mX-\mX')\\ S_u(v-v')\end{bmatrix}\\
&\quad+\begin{bmatrix}\Delta_qf\\ E_XC_x\Delta_qf\\ 0\end{bmatrix}
+\begin{bmatrix}\Delta_uf\\ E_XC_x\Delta_uf\\ E_u(\vu-\vu')\end{bmatrix}.
\end{aligned}
\end{equation*}
The squared norm of a stacked vector is the sum of the squared norms
of its blocks. The first vector therefore has norm at most
$\max\{\norm{S_X}_2,\norm{S_u}_2\}\norm{q-q'}_2$, the second at most
$\sqrt{1+\norm{E_XC_x}_2^2}\,\norm{\Delta_qf}_2$, and the third at most
$\big(b_B^2(1+\norm{E_XC_x}_2^2)+\norm{E_u}_2^2\big)^{1/2}
\norm{\vu-\vu'}_2$. The triangle inequality then gives
\begin{equation}
\begin{aligned}
&\norm{\mathcal F(q,\vu)-\mathcal F(q',\vu')}_2\\
&\quad\le L_q\norm{q-q'}_2+L_u\norm{\vu-\vu'}_2,
\end{aligned}
\label{eq:augmented_lipschitz}
\end{equation}
for all $q,q'\in\mathcal Q$ and $\vu,\vu'\in\mathcal U$, where
\begin{equation}
\begin{aligned}
L_q={}&\max\{\norm{S_X}_2,\norm{S_u}_2\}\\
&+\sqrt{1+\norm{E_XC_x}_2^2}\sqrt{a^2+\beta_C^2},\\
L_u={}&\sqrt{b_B^2(1+\norm{E_XC_x}_2^2)+\norm{E_u}_2^2}.
\end{aligned}
\label{eq:augmented_lipschitz_constants}
\end{equation}

\begin{lemma}[Freezing error bound]
\label{lem:freezing_error}
Suppose Assumptions~\ref{asm:domain}--\ref{asm:context_lipschitz} hold,
and let the frozen matrices $\bar\mA_{i|k}$ and $\bar\mB_{i|k}$ be
generated by the nominal rollout under $\bar U_k$ as
in~\eqref{eq:nominal_matrices}. Define
\begin{equation}
\begin{aligned}
K_i&=L_u\Big(\sum_{j=0}^{i-1}L_q^{2j}\Big)^{1/2},\\
D_0&=0,\qquad D_{i+1}=aD_i+\beta_C K_i.
\end{aligned}
\label{eq:freezing_constants}
\end{equation}
The empty sum gives $K_0=0$. Then, for every feasible candidate
sequence $U_k$ and nominal sequence $\bar U_k$,
\begin{equation}
\norm{\vz_{i|k}^P-\vz_{i|k}^F}_2\le D_i\delta_k,
\qquad\delta_k=\norm{U_k-\bar U_k}_2,
\label{eq:freezing_bound}
\end{equation}
for $i=0,\ldots,N$, and $D_1=0$.
\end{lemma}

\begin{proof}
The proof first bounds the deviation of the recursive trajectory from
the nominal one, and then propagates the resulting matrix mismatch
through the lifted error. Let $q_i^P=\col(\vz_i^P,\mX_i^P,v_i^P)$ and
$\bar q_i=\col(\bar\vz_i,\bar\mX_i,\bar v_i)$ be the augmented
recursive and nominal trajectories. By
Assumptions~\ref{asm:domain}--\ref{asm:bounds}, both trajectories
remain in $\mathcal Q$ for $i=0,\ldots,N$, so
\eqref{eq:augmented_lipschitz} applies at every step. Since
$q_0^P=\bar q_0$, iterating this estimate gives
\begin{equation}
\norm{q_i^P-\bar q_i}_2
\le L_u\sum_{j=0}^{i-1}L_q^{i-1-j}\norm{\vu_j-\bar\vu_j}_2
\le K_i\delta_k,
\label{eq:augmented_deviation}
\end{equation}
where the final inequality follows from Cauchy--Schwarz. We now turn
to the lifted error $\ve_i^F=\vz_i^P-\vz_i^F$. Subtracting the frozen
recursion~\eqref{eq:frozen_model} from the recursive one and
collecting terms yields
\begin{equation}
\begin{aligned}
\ve_{i+1}^F={}&\bar\mA_i\ve_i^F
+\big[\mA(\vc_i^P)-\bar\mA_i\big]\vz_i^P\\
&+\big[\mB(\vc_i^P)-\bar\mB_i\big]\vu_i.
\end{aligned}
\label{eq:freezing_recursion}
\end{equation}
The two bracketed terms are driven by the context deviation, which
satisfies
$\norm{\vc_i^P-\bar\vc_i}_2\le L_{Cq}\norm{q_i^P-\bar q_i}_2$.
Lemma~\ref{lem:scheduling_regularity} gives
$\norm{\bar\mA_i}_2\le a$ and bounds the two matrix differences by
$L_A\norm{\vc_i^P-\bar\vc_i}_2$ and $L_B\norm{\vc_i^P-\bar\vc_i}_2$.
With $\norm{\vz_i^P}_2\le M_P$, $\norm{\vu_i}_2\le M_u$, and
\eqref{eq:augmented_deviation}, it follows that
\begin{equation}
\norm{\ve_{i+1}^F}_2\le a\norm{\ve_i^F}_2+\beta_C K_i\delta_k.
\label{eq:freezing_scalar_recursion}
\end{equation}
Because $\ve_0^F=0$, induction on $i$
with~\eqref{eq:freezing_constants} gives~\eqref{eq:freezing_bound}.
Finally, $K_0=0$ and $D_0=0$ give $D_1=0$.
\end{proof}

The first-step freezing error vanishes because the three trajectories
share their initial histories and, by~\eqref{eq:past_input_history},
the context contains past inputs only. The first lifted update
therefore uses the same scheduled matrices under $U_k$ and
$\bar U_k$. Freezing errors at later prediction steps can nevertheless
affect the first optimized input through the horizon cost. The growth
of $D_i$ depends on both $a$ and $L_q$. Feasibility of $U_k$ and
$\bar U_k$ ensures neither a small $\delta_k$ nor that the predicted
histories remain in $\mathcal H$; the latter is the content of
Assumption~\ref{asm:domain}.

The frozen prediction is compared with the encoded physical
trajectory through the recursive predictor. The triangle inequality,
Theorem~\ref{thm:prediction_error}, and
Lemma~\ref{lem:freezing_error} give
\begin{equation}
\begin{aligned}
\norm{\vx_{i|k}^E-\vx_{i|k}^F}_2
&\le\norm{\vz_{i|k}^E-\vz_{i|k}^F}_2\\
&\le\norm{\vz_{i|k}^E-\vz_{i|k}^P}_2
+\norm{\vz_{i|k}^P-\vz_{i|k}^F}_2\\
&\le\eta_{i,k}(U_k):=\mathcal E_i+D_i\delta_k,
\end{aligned}
\label{eq:total_error}
\end{equation}
where $\vx_{i|k}^F=C_x\vz_{i|k}^F$. The term $\mathcal E_i$ accounts for the
closure residual of the recursive predictor, and $D_i\delta_k$ for the
scheduling freezing.

\subsection{Practical Closed-Loop Output Tracking Analysis}
\label{sec:closed_loop_stability}

We next relate the prediction mismatch to the MPC objective and
establish sufficient conditions for practical tracking of the physical
system. The argument has three steps. The prediction mismatch
in~\eqref{eq:total_error} is first converted into a uniform
perturbation of the MPC objective, which bounds the suboptimality of
the frozen optimal sequence on the physical system. A
cost-controllability condition on the physical system then turns this
suboptimality into a relaxed decrease of a finite-horizon value
function. Finally, the decrease yields finite-time and ultimate bounds
on the output tracking error. To distinguish physical trajectories
from the learned
predictor, introduce the true augmented state
$\xi_k=\col(\vs_k,\mX_k,v_k)$ and its transition map
\begin{equation}
\xi_{k+1}=\mathcal T(\xi_k,\vu_k)
=\begin{bmatrix}
f(\vs_k,\vu_k)\\
S_X\mX_k+E_Xh(f(\vs_k,\vu_k))\\
S_uv_k+E_u\vu_k
\end{bmatrix}.
\label{eq:true_augmented_transition}
\end{equation}
Unlike $\mathcal F$ in~\eqref{eq:augmented_update_map},
$\mathcal T$ describes the physical system and the measured history
updates. The physical state is introduced only for analysis and is
not required by Algorithm~\ref{alg:mpc}.

For the reference sequence $\{\vx_t^r\}_{t\ge0}$, assumed fixed across
successive preview windows and uniformly bounded by
$M_r=\sup_{t\ge0}\norm{\vx_t^r}_2<\infty$, define the ideal objective
\begin{equation}
\begin{aligned}
J_k^E(U;\xi_k)={}&\sum_{i=1}^N
\norm{C_x\vz_{i|k}^E(U)-\vx_{k+i}^r}_{Q_x}^2\\
&+\sum_{i=0}^{N-1}\norm{\vu_{i|k}}_{R_u}^2,
\end{aligned}
\label{eq:ideal_objective}
\end{equation}
where $U\in\mathcal U^N$ drives the true plant initialized at $\xi_k$.
The ideal and frozen objectives use the same cost and constraints,
but different prediction models. Let $J_k^F(U)$ denote the objective
in~\eqref{eq:quadratic_template} for a fixed feasible $\bar U_k$, with
the dependence on the available histories suppressed. The ideal
problem is a theoretical benchmark and is not solved online. Its
minimum, augmented by the current output cost, is the value function
of the closed-loop analysis; see~\eqref{eq:corrected_value_function}.

By compactness of the input set, define the uniform bounds
\begin{equation}
\bar\delta_N=\sqrt{N}\operatorname{diam}(\mathcal U),\qquad
\eta_i^{\max}=\mathcal E_i+D_i\bar\delta_N,
\label{eq:uniform_prediction_mismatch}
\end{equation}
where $\operatorname{diam}(\mathcal U)=\max_{\vu,\vu'\in\mathcal U}
\norm{\vu-\vu'}_2$. Since $\delta_k\le\bar\delta_N$, these bounds
cover every feasible candidate and nominal sequence. Moreover,
\eqref{eq:uniform_matrix_bound} gives the uniform frozen-state bound
$\norm{\vz_{i|k}^F(U)}_2\le M_F$, where
\begin{equation}
M_F=\max_{0\le i\le N}\left(a^iM_z+b_BM_u\sum_{j=0}^{i-1}a^j\right).
\label{eq:frozen_state_bound}
\end{equation}
This bound concerns the frozen lifted state only and is used to bound
the frozen output in the cost. It does not place the corresponding
histories in $\mathcal H$.

\begin{lemma}[Objective and optimal-input perturbation]
\label{lem:objective_policy_perturbation}
Suppose Assumptions~\ref{asm:domain}--\ref{asm:context_lipschitz}
hold, the physical transition and measurement maps $f$ and $h$
are continuous, and the reference is bounded by $M_r$. Consider any
initial condition and sampling instant covered by
Assumption~\ref{asm:domain}. Let $U_k^{E,*}$ be any minimizer of
$J_k^E$ and $U_k^{F,*}$ the unique minimizer of $J_k^F$ on the
nonempty, compact, and convex set $\mathcal U^N$.
Define
\begin{equation}
\bar\Delta_N=\lambda_{\max}(Q_x)\sum_{i=1}^N
\Big[2(M_F+M_r)\eta_i^{\max}+(\eta_i^{\max})^2\Big].
\label{eq:uniform_cost_perturbation}
\end{equation}
For every feasible $U$,
\begin{equation}
|J_k^E(U)-J_k^F(U)|\le\bar\Delta_N.
\label{eq:cost_perturbation_sum}
\end{equation}
Furthermore,
\begin{gather}
0\le J_k^E(U_k^{F,*})-J_k^E(U_k^{E,*})\le2\bar\Delta_N,
\label{eq:ideal_suboptimality}\\
\norm{U_k^{E,*}-U_k^{F,*}}_2\le2\sqrt{\bar\Delta_N/\mu},
\qquad \mu=2\lambda_{\min}(R_u),
\label{eq:optimal_input_perturbation}
\end{gather}
where $\lambda_{\max}$ and $\lambda_{\min}$ denote the largest and
smallest eigenvalues of the symmetric weighting matrices. In
particular, the first input blocks satisfy
\begin{equation}
\norm{\vu_{0|k}^{E,*}-\vu_{0|k}^{F,*}}_2
\le2\sqrt{\bar\Delta_N/\mu}.
\label{eq:first_input_perturbation}
\end{equation}
\end{lemma}

\begin{proof}
The proof bounds the objective difference term by term, converts it
into a suboptimality gap, and then uses strong convexity to bound the
distance between the minimizers. For each $i$, let
$\vp_i=C_x\vz_{i|k}^F(U)-\vx_{k+i}^r$ and
$\ve_i=C_x(\vz_{i|k}^E(U)-\vz_{i|k}^F(U))$. By~\eqref{eq:total_error}
and~\eqref{eq:frozen_state_bound}, $\norm{\ve_i}_2\le\eta_i^{\max}$
and $\norm{\vp_i}_2\le M_F+M_r$. Since the input penalties coincide,
the two objectives differ only in their output costs. Expanding each
square gives
\begin{equation}
\begin{aligned}
&|\norm{\vp_i+\ve_i}_{Q_x}^2-\norm{\vp_i}_{Q_x}^2|\\
&\quad\le\lambda_{\max}(Q_x)\big(2\norm{\vp_i}_2\norm{\ve_i}_2
+\norm{\ve_i}_2^2\big).
\end{aligned}
\label{eq:cost_perturbation}
\end{equation}
Summing over $i=1,\ldots,N$ proves~\eqref{eq:cost_perturbation_sum}.

For the suboptimality gap, note first that the ideal objective is
continuous in $U$, by the continuity of $f$ and $h$ and the
state-preserving readout. A minimizer $U_k^{E,*}$ therefore exists on
the compact set $\mathcal U^N$, and no convexity of the ideal problem
is required. The optimality of $U_k^{F,*}$ for $J_k^F$,
with~\eqref{eq:cost_perturbation_sum} applied twice, yields
\begin{equation*}
\begin{aligned}
J_k^E(U_k^{F,*})&\le J_k^F(U_k^{F,*})+\bar\Delta_N
\le J_k^F(U_k^{E,*})+\bar\Delta_N\\
&\le J_k^E(U_k^{E,*})+2\bar\Delta_N.
\end{aligned}
\end{equation*}
The lower bound in~\eqref{eq:ideal_suboptimality} is the optimality
of $U_k^{E,*}$ for $J_k^E$.

It remains to bound the distance between the two minimizers. With
fixed matrices, the frozen prediction is affine in $U$ and the
input penalty contributes $2(I_N\otimes R_u)$ to the Hessian, where
$I_N$ is the $N$-dimensional identity and $\otimes$ is the Kronecker
product. Thus $J_k^F$ is
$\mu$-strongly convex. Because $\mathcal U^N$ is convex, its
constrained minimizer satisfies the first-order variational inequality
$\nabla J_k^F(U_k^{F,*})^\top(U-U_k^{F,*})\ge0$ for all
$U\in\mathcal U^N$, rather than a vanishing gradient. Combined with
strong convexity, this inequality gives
\begin{equation}
\begin{aligned}
&\frac{\mu}{2}\norm{U_k^{E,*}-U_k^{F,*}}_2^2
\le J_k^F(U_k^{E,*})-J_k^F(U_k^{F,*})\\
&\quad\le J_k^E(U_k^{E,*})-J_k^E(U_k^{F,*})+2\bar\Delta_N
\le2\bar\Delta_N.
\end{aligned}
\label{eq:strong_convexity_gap}
\end{equation}
Here the second inequality applies~\eqref{eq:cost_perturbation_sum} to
both objective values, and the third uses the optimality of
$U_k^{E,*}$ for $J_k^E$. This
proves~\eqref{eq:optimal_input_perturbation}. Since selecting the
first input block cannot increase the Euclidean norm,
\eqref{eq:first_input_perturbation} follows.
\end{proof}

Lemma~\ref{lem:objective_policy_perturbation} connects the prediction
error with both the suboptimality of the realized input sequence on
the physical system and the deviation from an ideal optimal sequence.
We use the suboptimality bound directly in the closed-loop analysis,
which avoids requiring Lipschitz continuity of the ideal optimal
control law or of a Lyapunov function composed with the plant.

Let $p_t(\xi)=\norm{h(\vs)-\vx_t^r}_{Q_x}^2$ denote the current
output cost for $\xi=\col(\vs,\mX,v)$, and define the stage cost
$\ell_t(\xi,\vu)=p_t(\xi)+\norm{\vu}_{R_u}^2$. For a time $t$ and
a horizon $j\ge1$, define the true-system value function
\begin{equation}
W_{j,t}(\xi)=\min_{U\in\mathcal U^j}\bigg[\sum_{i=0}^{j-1}
\ell_{t+i}(\xi_{i|t},\vu_{i|t})+p_{t+j}(\xi_{j|t})\bigg],
\label{eq:ideal_value_function}
\end{equation}
where $\xi_{0|t}=\xi$ and
$\xi_{i+1|t}=\mathcal T(\xi_{i|t},\vu_{i|t})$. We also set
$W_{0,t}=p_t$. For horizon $N$, write $W_k=W_{N,k}$. In particular,
\begin{equation}
W_k(\xi)=p_k(\xi)+\min_{U\in\mathcal U^N}J_k^E(U;\xi).
\label{eq:corrected_value_function}
\end{equation}
The added current output cost is independent of the decision variables,
so this identity does not modify the MPC objective or its optimizer.
It gives the lower bound
\begin{equation}
W_k(\xi)\ge p_k(\xi)\ge\lambda_{\min}(Q_x)
\norm{h(\vs)-\vx_k^r}_2^2.
\label{eq:lyapunov_output_lower_bound}
\end{equation}

For MPC without a stabilizing terminal constraint, a sufficiently long
horizon may yield a relaxed decrease condition under appropriate
cost-controllability hypotheses~\cite{grune2010analysis,coron2020model}.
For the present tracking objective, we use the following finite-horizon
condition with an additive allowance for nonvanishing tracking costs.

\begin{assumption}[Cost controllability with an additive allowance]
\label{asm:cost_controllability}
Let $N\ge2$ and let $\mathcal O$ be a closed operating domain of
consistent true augmented states. There exist $B_{\rm cc}>1$ and
$d_{\rm cc}\ge0$ such that
\begin{equation}
W_{j,t}(\xi)\le B_{\rm cc}\min_{\vu\in\mathcal U}\ell_t(\xi,\vu)
+d_{\rm cc},\qquad j=1,\ldots,N,
\label{eq:cost_controllability}
\end{equation}
for all $t\ge0$ and $\xi\in\mathcal O$, and at the intermediate
states reached at steps $i=1,\ldots,N-1$ by any feasible
$N$-step input sequence initialized in $\mathcal O$ (with $t$
replaced by the corresponding time). The constants are uniform over
these states and times.
\end{assumption}

This assumption concerns the physical system, reference sequence,
input constraints, and cost. In particular, it is not implied by
matrix boundedness or controllability of the lifted predictor. The
allowance $d_{\rm cc}$ accommodates nonzero reference-maintaining
inputs and discrepancies not characterized by the current measurement
error. For example, a sufficient way to verify
\eqref{eq:cost_controllability} is to construct a feasible comparison
trajectory for each such initial state and time. If, for constants
$C_0>0$, $0\le\sigma<1$, and $\bar\ell\ge0$, its stage costs satisfy
$\ell_{t+i}(\xi_{i|t},\vu_{i|t})\le C_0\sigma^i p_t(\xi)+\bar\ell$
and its terminal output cost satisfies the same bound at $i=j$,
then summation gives~\eqref{eq:cost_controllability} with any
$B_{\rm cc}\ge C_0/(1-\sigma)$, $B_{\rm cc}>1$, and
$d_{\rm cc}=(N+1)\bar\ell$. This is a sufficient comparison condition,
not a property asserted for every partially observed plant. In a
specific application, the constants must hold uniformly over all
initial states, times, and intermediate states named in
Assumption~\ref{asm:cost_controllability} for the chosen dynamic
reference; a finite collection of closed-loop trajectories does not
by itself establish them.

\begin{assumption}[Closed-loop domain]
\label{asm:closed_loop_domain}
The conditions in Assumptions~\ref{asm:domain}--\ref{asm:context_lipschitz} hold
uniformly for initial true augmented states in $\mathcal O$ and all
feasible candidate and nominal input sequences. For some $R>0$,
let $\Omega_{R,k}=\{\xi\in\mathcal O:W_k(\xi)\le R\}$. Assume
$\xi_0\in\Omega_{R,0}$ and
$\mathcal T(\xi,\vu)\in\mathcal O$ for every
$\xi\in\Omega_{R,k}$, $\vu\in\mathcal U$, and $k\ge0$.
\end{assumption}

The successor-domain condition is a local sufficient hypothesis
ensuring that the estimates can be applied along the physical closed
loop. It is assumed rather than derived from the controller, and it is
distinct from feasibility of the QP, which follows from the input-only
constraints. Given this condition, the theorem below establishes
forward invariance of the sublevel sets $\Omega_{R,k}$.

\begin{theorem}[Practical closed-loop output tracking]
\label{thm:practical_stability}
Suppose Assumptions~\ref{asm:domain}--\ref{asm:closed_loop_domain}
and the continuity and bounded-reference conditions of
Lemma~\ref{lem:objective_policy_perturbation} hold. Define
\begin{equation}
\begin{aligned}
\vartheta&=1-1/B_{\rm cc},\\
\lambda_N&=1/B_{\rm cc}-(B_{\rm cc}-1)\vartheta^{N-1},\\
d_N&=\big[1+1/B_{\rm cc}
+(B_{\rm cc}-1)(1-\vartheta^{N-1})\big]d_{\rm cc}\\
&\quad+2B_{\rm cc}\bar\Delta_N.
\end{aligned}
\label{eq:robust_certificate_constants}
\end{equation}
If $\lambda_N>0$ and $d_N\le\lambda_NR$, then the physical closed
loop under Algorithm~\ref{alg:mpc} satisfies
\begin{equation}
W_{k+1}(\xi_{k+1})\le(1-\lambda_N)W_k(\xi_k)+d_N.
\label{eq:robust_lyapunov_decrease}
\end{equation}
The time-varying sublevel sets $\Omega_{R,k}$ are forward invariant
under the controller for every feasible initial nominal sequence, and
\begin{equation}
\begin{aligned}
\norm{\vx_k-\vx_k^r}_2^2
\le{}&\frac{1}{\lambda_{\min}(Q_x)}\Big[ (1-\lambda_N)^kW_0(\xi_0)\\
&+\frac{d_N}{\lambda_N}\big(1-(1-\lambda_N)^k\big)\Big].
\end{aligned}
\label{eq:finite_time_tracking_bound}
\end{equation}
In particular,
\begin{equation}
\limsup_{k\to\infty}\norm{\vx_k-\vx_k^r}_2
\le\sqrt{\frac{d_N}{\lambda_N\lambda_{\min}(Q_x)}}.
\label{eq:conditional_ultimate_bound}
\end{equation}
\end{theorem}

\begin{proof}
The proof follows the true trajectory driven by the frozen optimal
sequence. It first shows that every suffix of this sequence is nearly
optimal for its remaining horizon, then uses cost controllability to
bound the last suffix, and finally constructs a comparison sequence at
the successor state. At time $k$, let $U_k^{F,*}$ be the frozen optimal
sequence and let $\widetilde\xi_i=\xi_{i|k}^{E}$ denote the true
augmented trajectory driven by that sequence from
$\widetilde\xi_0=\xi_k$. Write
$\ell_i=\ell_{k+i}(\widetilde\xi_i,\vu_{i|k}^{F,*})$ for the stage
costs along this trajectory, and define the remaining cost at step $i$
by
\begin{equation}
S_i=\sum_{j=i}^{N-1}\ell_j+p_{k+N}(\widetilde\xi_N),
\qquad i=0,\ldots,N.
\label{eq:remaining_cost}
\end{equation}
The scalar $S_i$ is unrelated to the history shift matrices $S_X$ and
$S_u$. By~\eqref{eq:corrected_value_function},
$S_0=p_k(\xi_k)+J_k^E(U_k^{F,*})$, so
\eqref{eq:ideal_suboptimality} gives, with $\omega_N=2\bar\Delta_N$,
\begin{equation}
W_k(\xi_k)\le S_0\le W_k(\xi_k)+\omega_N.
\label{eq:initial_value_gap}
\end{equation}
The same margin holds for every suffix. Indeed, retaining the first
$i$ inputs and replacing the remaining ones by an ideal optimizer
from $\widetilde\xi_i$ gives a feasible sequence of length $N$ for the
true problem. Hence
\begin{equation*}
W_k(\xi_k)\le S_0-S_i+W_{N-i,k+i}(\widetilde\xi_i),
\end{equation*}
which, together with~\eqref{eq:initial_value_gap}, implies
\begin{equation}
S_i\le W_{N-i,k+i}(\widetilde\xi_i)+\omega_N,
\qquad i=0,\ldots,N-1.
\label{eq:suffix_suboptimality}
\end{equation}
Assumption~\ref{asm:cost_controllability} applies at $\xi_k$ and at
the intermediate states $\widetilde\xi_i$. Because
$\min_{\vu\in\mathcal U}\ell_{k+i}(\widetilde\xi_i,\vu)\le\ell_i$,
it bounds the value function in~\eqref{eq:suffix_suboptimality} and
gives $S_i\le B_{\rm cc}\ell_i+d_{\rm cc}+\omega_N$. Since
$S_{i+1}=S_i-\ell_i$, this yields the recursion
\begin{equation*}
S_{i+1}\le\vartheta S_i+(d_{\rm cc}+\omega_N)/B_{\rm cc}.
\end{equation*}
Iterating from $i=0$ to $N-2$ and using~\eqref{eq:initial_value_gap}
yields
\begin{equation}
S_{N-1}\le\vartheta^{N-1}W_k(\xi_k)
+(1-\vartheta^{N-1})d_{\rm cc}+\omega_N.
\label{eq:last_suffix_bound}
\end{equation}

We now turn to the successor state. Under the receding-horizon
implementation, $\xi_{k+1}=\widetilde\xi_1$ is the actual successor.
To upper-bound the value function there, retain the inputs
$\vu_{1|k}^{F,*},\ldots,\vu_{N-2|k}^{F,*}$ and complete the sequence
with a two-step ideal optimizer from $\widetilde\xi_{N-1}$ at time
$k+N-1$; for $N=2$, the retained prefix is empty. Because the problem
has input constraints only, this comparison sequence is feasible. It
serves the proof alone and does not replace the nominal shift
in~\eqref{eq:nominal_shift}. Therefore,
\begin{equation}
\begin{aligned}
W_{k+1}(\xi_{k+1})
&\le S_0-\ell_0-S_{N-1}+W_{2,k+N-1}(\widetilde\xi_{N-1})\\
&\le S_0-\ell_0+(B_{\rm cc}-1)S_{N-1}+d_{\rm cc},
\end{aligned}
\label{eq:ideal_shift_identity}
\end{equation}
where the second inequality follows from
$W_{2,k+N-1}(\widetilde\xi_{N-1})\le B_{\rm cc}\ell_{N-1}+d_{\rm cc}$ and
$\ell_{N-1}\le S_{N-1}$. It remains to bound $\ell_0$ from below.
Cost controllability at the initial state gives
$W_k(\xi_k)\le B_{\rm cc}\ell_0+d_{\rm cc}$, that is,
$\ell_0\ge(W_k(\xi_k)-d_{\rm cc})/B_{\rm cc}$. Substituting this
estimate, together with~\eqref{eq:initial_value_gap} and
\eqref{eq:last_suffix_bound}, into~\eqref{eq:ideal_shift_identity}
results in
\begin{equation*}
\begin{aligned}
&W_{k+1}(\xi_{k+1})
\le\big[1-1/B_{\rm cc}+(B_{\rm cc}-1)\vartheta^{N-1}\big]W_k(\xi_k)\\
&\quad+\big[1+1/B_{\rm cc}
+(B_{\rm cc}-1)(1-\vartheta^{N-1})\big]d_{\rm cc}+B_{\rm cc}\omega_N.
\end{aligned}
\end{equation*}
This is~\eqref{eq:robust_lyapunov_decrease} with the constants
in~\eqref{eq:robust_certificate_constants}, since
$B_{\rm cc}\omega_N=2B_{\rm cc}\bar\Delta_N$.

The decrease inequality was derived for $\xi_k\in\mathcal O$, so it
must be shown to persist along the closed loop. For
$\xi_k\in\Omega_{R,k}$, Assumption~\ref{asm:closed_loop_domain}
ensures $\xi_{k+1}\in\mathcal O$. Since $0<\lambda_N<1$ and
$d_N\le\lambda_NR$,~\eqref{eq:robust_lyapunov_decrease} gives
$W_{k+1}(\xi_{k+1})\le(1-\lambda_N)R+\lambda_NR=R$, that is,
$\xi_{k+1}\in\Omega_{R,k+1}$. The nominal sequence also remains
feasible by~\eqref{eq:nominal_shift}. Induction on $k$ therefore
establishes forward invariance and validates all uniform estimates at
subsequent sampling instants. Iterating the scalar decrease inequality
and using~\eqref{eq:lyapunov_output_lower_bound} proves
\eqref{eq:finite_time_tracking_bound}, and letting $k\to\infty$ gives
\eqref{eq:conditional_ultimate_bound}.
\end{proof}

The two conditions of Theorem~\ref{thm:practical_stability} play
different roles. The condition $\lambda_N>0$ makes
\eqref{eq:robust_lyapunov_decrease} a contraction of the value function
up to the offset $d_N$, whereas $d_N\le\lambda_NR$ keeps
$W_{k+1}(\xi_{k+1})\le R$ and hence the successor state inside the
certified sublevel set. The horizon condition $\lambda_N>0$ is
equivalent to
$B_{\rm cc}(B_{\rm cc}-1)\vartheta^{N-1}<1$. For a fixed finite
cost-controllability constant it holds for sufficiently large $N$.
However, enlarging $N$ can also enlarge $\bar\Delta_N$ and
$d_{\rm cc}$, and the domain condition remains necessary. Thus a longer
horizon does not automatically yield a smaller practical tracking
bound. Theorem~\ref{thm:practical_stability} relates the final tracking
bound explicitly to the modeling and freezing errors through
$\bar\Delta_N$ and to the cost-controllability allowance through
$d_{\rm cc}$. If both $d_{\rm cc}$ and $\bar\Delta_N$ vanish under the stated
hypotheses, the output tracking error converges to zero. When $0\in\mathcal U$, Assumption~\ref{asm:cost_controllability}
also gives $W_k(\xi)\le B_{\rm cc}p_k(\xi)+d_{\rm cc}$. Hence, in the
zero-disturbance case,~\eqref{eq:finite_time_tracking_bound} establishes
exponential output tracking within the certified sublevel sets. For
nonzero allowances, the result is ultimate output tracking boundedness,
not asymptotic stability of the unmeasured physical state. Related data-driven MPC stability results
also rely on additional hypotheses regarding error bounds and cost
controllability~\cite{schimperna2026data,liu2026lpv}.

In particular, the absolute input penalty in~\eqref{eq:quadratic_template}
can produce a tracking offset when the reference requires nonzero
inputs. Also, an arbitrary reference change beyond the current preview
can increase the cost of any repeated-terminal-input completion. The
analysis therefore does not infer a terminal decrease from the
feasibility of the nominal shift, and it does not guarantee zero
tracking error for an arbitrary bounded reference. Verifying the
cost-controllability and operating-domain conditions for the target
plant and reference is necessary for applying the closed-loop bound.

\endgroup

\begin{thebibliography}{99}
\bibitem{koopman1931hamiltonian}
B. O. Koopman, ``Hamiltonian systems and transformation in Hilbert space,'' \emph{Proc. Natl. Acad. Sci. USA}, vol. 17, no. 5, pp. 315--318, 1931.
\bibitem{mezic2005spectral}
I. Mezi\'c, ``Spectral properties of dynamical systems, model reduction and decompositions,'' \emph{Nonlinear Dyn.}, vol. 41, no. 1--3, pp. 309--325, 2005.
\bibitem{brunton2022modern}
S. L. Brunton, M. Budi\v{s}i\'c, E. Kaiser, and J. N. Kutz, ``Modern Koopman theory for dynamical systems,'' \emph{SIAM Rev.}, vol. 64, no. 2, pp. 229--340, 2022.
\bibitem{williams2015data}
M. O. Williams, I. G. Kevrekidis, and C. W. Rowley, ``A data-driven approximation of the Koopman operator: Extending dynamic mode decomposition,'' \emph{J. Nonlinear Sci.}, vol. 25, no. 6, pp. 1307--1346, 2015.
\bibitem{proctor2016dynamic}
J. L. Proctor, S. L. Brunton, and J. N. Kutz, ``Dynamic mode decomposition with control,'' \emph{SIAM J. Appl. Dyn. Syst.}, vol. 15, no. 1, pp. 142--161, 2016.
\bibitem{vaswani2017attention}
A. Vaswani \emph{et al.}, ``Attention is all you need,'' in \emph{Proc. NeurIPS}, 2017, pp. 5998--6008.
\bibitem{grune2010analysis}
L. Gr{\"u}ne, J. Pannek, M. Seehafer, and K. Worthmann, ``Analysis of unconstrained nonlinear MPC schemes with time varying control horizon,'' \emph{SIAM J. Control Optim.}, vol. 48, no. 8, pp. 4938--4962, 2010.
\bibitem{coron2020model}
J.-M. Coron, L. Gr{\"u}ne, and K. Worthmann, ``Model predictive control, cost controllability, and homogeneity,'' \emph{SIAM J. Control Optim.}, vol. 58, no. 5, pp. 2979--2996, 2020.
\bibitem{schimperna2026data}
I. Schimperna, K. Worthmann, M. Schaller, L. Bold, and L. Magni, ``Data-driven model predictive control: Asymptotic stability despite approximation errors exemplified in the Koopman framework,'' \emph{Automatica}, vol. 193, art. 113203, 2026.
\bibitem{liu2026lpv}
Z. Liu, W. Li, J. Song, and Y. Liu, ``LPV-Koopman model predictive control: Modeling error and stability analysis,'' \emph{J. Franklin Inst.}, vol. 363, no. 12, art. 108791, 2026.
\end{thebibliography}
\end{document}
